\documentclass[a4paper]{article}     
\usepackage{amsfonts}
\usepackage{amsmath}
\usepackage{amssymb}
\usepackage{graphicx}

\begin{document}

\noindent \large Auteur: Mateusz Ragan \\
\noindent \large Studierichting: Informatica\\
\noindent \large Oefeningenreeks 1D nr. 24.2\\

\medskip

\normalsize

Los op in $\mathbb{C}$\\

$8z^6 + 7iz^3 + 1 = 0$  stel $t = z^3$\\

$8t^2 + 7it + 1 = 0$\\

$\Delta = (7i)^2 - 4\cdot 8 \cdot 1 = -81$\\

$t_1 = \dfrac{-7i + 9i}{16} = \dfrac{i}{8}$ \\

$t_2 = \dfrac{-7i - 9i}{16} = -i$\\

$8\left(t-\dfrac{i}{8}\right)(t + i) = 0 \Leftrightarrow (8t-i)(t+i) = 0$\\

In het begin stelden we $t = z^3$ dus nu moeten we 2 derdemachtswortels oplossen. \\

$(8z^3 - i)(z^3 + 1) = 0$\\

We beginnen eerst met de eerste vergelijking.\\

$8z^3 - i = 0 \Leftrightarrow z^3 = \dfrac{i}{8}$\\

$z = \dfrac{1}{2}$ cis$\left(\dfrac{\dfrac{\pi}{2}+k\cdot 2\pi}{3}\right)$ met $k \in \{0,1,2\}$\\

$w_0 = \dfrac{1}{2}$ cis$\left(\dfrac{\pi}{6}\right) = \dfrac{\sqrt{3}}{4} + \dfrac{i}{4}$ \\

$w_1 = \dfrac{1}{2}$ cis$\left(\dfrac{5\pi}{6}\right) = - \dfrac{\sqrt{3}}{4} + \dfrac{i}{4}$ \\

$w_2 = \dfrac{1}{2}$ cis$\left(\dfrac{3\pi}{2}\right) = - \dfrac{i}{2}$ \\

Dus voor deze vergelijking zijn er 3 oplossingen.\\

$z = \dfrac{\sqrt{3}}{4} + \dfrac{i}{4}  \vee z = - \dfrac{\sqrt{3}}{4} + \dfrac{i}{4} \vee z = - \dfrac{i}{2}$\\

Nu lossen we de volgende vergelijking op het volgende blad.\\

\newpage

$z^3+i = 0 \Leftrightarrow z^3 = -i$\\

$z$ = cis$\left(\dfrac{- \dfrac{\pi}{2}+k\cdot 2\pi}{3}\right)$ met $k \in \{0,1,2\}$\\

$w_0 = $ cis$\left(\dfrac{- \pi}{6}\right) = \dfrac{\sqrt{3}}{2} - \dfrac{i}{2}$ \\

$w_1 = $ cis$\left(\dfrac{\pi}{2}\right) = i $ \\

$w_2 = $ cis$\left(\dfrac{7\pi}{6}\right) = - \dfrac{\sqrt{3}}{2} - \dfrac{i}{2}$ \\

Dus voor deze vergelijking zijn er 3 oplossingen.\\

$z = \dfrac{\sqrt{3}}{2} - \dfrac{i}{2}   \vee z =  i  \vee z = - \dfrac{\sqrt{3}}{2} - \dfrac{i}{2}$\\

Dus voor de vergelijking: $8z^6 + 7iz^3 + 1 = 0$ zijn er 6 oplossingen.\\

$$z = \dfrac{\sqrt{3}}{4} + \dfrac{i}{4} \vee z = - \dfrac{\sqrt{3}}{4} + \dfrac{i}{4} \vee z = - \dfrac{i}{2} \vee z =  \dfrac{\sqrt{3}}{2} - \dfrac{i}{2} \vee z = i \vee z = - \dfrac{\sqrt{3}}{2} - \dfrac{i}{2}$$

\begin{thebibliography}{999}
%\bibitem{a} Referentie
\end{thebibliography}
\end{document} 
